\chapter{Introduzione}
\label{chap:introduction}

\paragraph{Riassunto del capitolo.} Questo capitolo inquadra il
problema: la caratterizzazione di celle Litio-Ione tramite spettroscopia
di impedenza è costosa in termini di tempo macchina e operatore, e
genera dataset alto-dimensionali (Aging, Temperatura, SOC, Frequenza)
in cui non sempre tutte le combinazioni sono disponibili. Il progetto,
svolto in collaborazione col Politecnico di Milano, parte da questa
esigenza e definisce tre obiettivi predittivi: ricostruzione di curve
mancanti, classificazione binaria young/old, interpolazione di un
intero livello di invecchiamento. Sopra al deliverable formale (tre
notebook Jupyter) abbiamo poi costruito \emph{di nostra iniziativa}
una piccola piattaforma MLOps end-to-end, oggetto del
\cref{chap:mlops}.

\section{Contesto}

Le batterie Litio-Ione sono il sostrato energetico dominante della
mobilità elettrica, dei dispositivi consumer e dello stoccaggio
stazionario. La loro caratterizzazione lungo il ciclo di vita richiede
strumenti capaci di catturare \emph{processi elettrochimici eterogenei}
(trasporto di carica, doppio strato, diffusione) che si manifestano su
scale di frequenza differenti. Lo strumento più informativo e meno
invasivo è la \textbf{Galvanostatic Electrochemical Impedance
Spectroscopy (\acrshort{geis})}: si inietta una piccola corrente
sinusoidale a molte frequenze e si misura la risposta in tensione,
ricavando l'impedenza complessa
$Z(f) = \mathrm{Re}(Z) + j\,\mathrm{Im}(Z)$. Il piano di Nyquist
$(\mathrm{Re}(Z), -\mathrm{Im}(Z))$ traduce queste misure in una firma
geometrica che evolve con \emph{aging}, \emph{temperatura} e \emph{state
of charge}.

Acquisire un dataset GEIS è costoso: ogni curva richiede minuti di
strumentazione di precisione, pilotaggio termico stabile e protezioni
della cella. Una caratterizzazione completa di una pouch
\acrshort{licoo2} su una matrice $5 \times 8 \times 5$
(aging $\times$ temperatura $\times$ SOC) impiega nell'ordine delle
decine di ore. Diventa quindi rilevante chiedersi:
\begin{itemize}
  \item \emph{Possiamo evitare di misurare alcune coppie (Aging,
        Temperatura) e ricostruirle dai dati esistenti?}
  \item \emph{Dati un livello di Aging e un livello di SOC mai
        osservati in addestramento, possiamo classificare come
        ``young'' o ``old'' le 8 curve di Nyquist (una per temperatura)
        di quella coppia, leggendo solo la loro forma?}
  \item \emph{Possiamo prevedere un intero stadio di invecchiamento mai
        osservato direttamente?}
\end{itemize}

Il dataset utilizzato in questo progetto è stato fornito dal
\textbf{Politecnico di Milano} nell'ambito di una collaborazione di
ricerca; il progetto è stato sviluppato presso \textbf{SUPSI -- DTI}
come progetto di semestre.

\section{Obiettivi del progetto}

\subsection{Compito primario}

Il compito primario era \textbf{sviluppare tre notebook Jupyter} -- uno
per ciascuna delle tre domande sopra elencate -- in cui addestrare e
valutare modelli di machine learning capaci di:
\begin{enumerate}[leftmargin=*]
  \item ricostruire una coppia $(Aging, Temperatura)$ esclusa dal
        training (regressione multi-output, Leave-One-Out);
  \item dati un livello di Aging e un livello di SOC, classificare
        come \emph{young} (Aging $\leq 2$) o \emph{old}
        (Aging $\geq 3$) le \textbf{8 curve di Nyquist} (una per
        temperatura) di quella coppia, tenute fuori dall'addestramento
        (classificazione binaria);
  \item interpolare un \emph{intero} livello di invecchiamento mai
        osservato in addestramento (regressione multi-output, Leave-One-Aging-Out).
\end{enumerate}

Per ogni task il deliverable era un benchmark di modelli con metriche
d'accuratezza, scelta del modello vincente ed una motivazione tecnica.

\subsection{Estensione volontaria (non richiesta)}

Una volta consolidati i tre notebook, abbiamo deciso \emph{di nostra
iniziativa} di costruire intorno a essi una piccola piattaforma per
mostrare i risultati a un utente non-developer e per esercitarci nelle
prassi MLOps di settore. Questa estensione comprende:
\begin{itemize}
  \item un package Python (\codefile{src/}) che incapsula la logica dei
        notebook in funzioni testabili e riutilizzabili;
  \item un backend HTTP \texttt{FastAPI}~\cite{fastapi} con OpenAPI e CORS
        configurabile;
  \item un frontend \texttt{React + TypeScript + Vite} che consuma le
        \acrshort{api} e renderizza grafici Plotly;
  \item experiment tracking con MLflow~\cite{mlflow2018} (parent run per
        task, child run per modello);
  \item suite di 81 test pytest con coverage gate al 90\,\% in CI (95\,\% attuale);
  \item containerizzazione (Dockerfile multistage + compose);
  \item un modulo di \textbf{drift detection} (KS + PSI) con CLI ed
        endpoint REST.
\end{itemize}

Va sottolineato che \textbf{questo livello di ingegnerizzazione non era
richiesto}; la consegna formale era costituita dai soli notebook. Le
sezioni che seguono trattano insieme entrambe le parti: il \emph{cuore}
ML predittivo e l'\emph{infrastruttura} costruita sopra.
